\section{Introduction}
\label{sec:intro}

Removing transformer blocks and then continuing pretraining (CPT, often called
``healing'') is an attractive route to smaller language models
~\citep{gromov2024unreasonable,shortenedllama,men2024shortgpt,
yang2024laco}. A pruned model can improve rapidly on held-out next-token loss,
and endpoint perplexity is therefore a natural measure of whether retraining is
working. It is not, however, a complete description of what the structural
intervention removed or what CPT restored.

Perplexity averages local prediction loss over a corpus. Factual recall,
multiple-choice answer selection, generation, and task-specific interfaces
probe different observables. Compression studies have documented that
aggregate language-model metrics can miss changes in parametric knowledge,
downstream behavior, or safety~\citep{costcompression,beyondperplexity,
jaiswal2024truth}. Yet most depth-pruning work reports a small number of
endpoints. Endpoint comparisons cannot distinguish slow recovery from a stable
deficit, corpus drift from structural damage, or competence loss from failure
to express an answer through a particular scoring interface.

\paragraph{Research question.}
We ask: \emph{what does continued pretraining restore after depth pruning, at
what observed rate, and which controls are required to interpret the result?}
We treat recovery as a multi-axis object---held-out likelihood, target
capability, scoring interface, structural policy, and training budget---rather
than a scalar certified by perplexity.

\begin{figure*}[t]
\centering
\IfFileExists{figures/fig1_pruneheal_autofigure.pdf}
  {\includegraphics[width=\textwidth]{figures/fig1_pruneheal_autofigure.pdf}}
  {\includegraphics[width=\textwidth]{figures/fig1_concept_gptimage_corrected.pdf}}
\caption{\textbf{Post-pruning recovery is multi-axis.} The principal recipe
keeps 14 of 32 pretrained OLMo-2 blocks, appends two fresh blocks, and continues
pretraining. At 200k steps, held-out PPL improves sharply but remains
$1.428\times$ the base, while answer-letter MMLU recovers only 19.5\% of the
base's above-chance margin. A different 16-layer construction, ShortGPT,
reaches 63.0\%, showing that nominal depth does not determine one endpoint.
Layer-wise readout markers are descriptive
correlates, not causal storage boundaries.}
\label{fig:concept}
\end{figure*}

We instantiate the study on OLMo-2-1124-7B~\citep{olmo2}. Prefix arms retain
$k$ pretrained decoder blocks, append two randomly initialized blocks, and
continue pretraining on a fixed Dolmino stream. The principal keep14 arm is
tracked to 200k steps. Controls include the intact base, an unpruned model
continued on the same corpus array and nominal schedule to a pre-registered
25k plateau checkpoint, a
frozen-prefix arm, a fully random same-shape model, and a non-contiguous
ShortGPT-16 endpoint. Evaluation combines held-out likelihood, a broad zero-shot
likelihood suite, answer-letter and complete-option MMLU, and no-retrieval
PopQA, TriviaQA, and NQ-open.

The central observation is bounded but robust across interfaces. keep14 reaches
PPL 10.561, yet remains far below the intact model on answer-letter MMLU and
all three closed-book tasks. Complete-option scoring narrows the MMLU gap,
consistent with an answer-symbol or readout component, but the random-init
content-score floor rules out interpreting that protocol as recovered subject
knowledge. ShortGPT is substantially stronger on both PPL and MMLU, so the
keep14 deficit is not a universal consequence of using 16 blocks. Coupled
structural differences prevent a single-factor explanation.

\paragraph{Contributions.}
\begin{itemize}[leftmargin=1.2em,itemsep=2pt]
\item \textbf{Trajectory-level measurement.} We track likelihood and
knowledge-sensitive behavior during post-depth-pruning CPT rather than
evaluating only a final compressed endpoint.
\item \textbf{Controls that separate interpretations.} Continued-full-model,
same-shape initialization/adaptation, scoring-interface, policy, and
closed-book controls distinguish corpus drift, fluency floors, answer format,
and structural choice.
\item \textbf{A reporting protocol.} We show why post-pruning likelihood,
target capability, interface, exact structure, and recovery compute should be
reported as separate axes. Layer-wise readouts provide a compatible but
explicitly non-causal correlate.
\end{itemize}
